% !TEX program = xelatex
\documentclass{kawcv}

\usepackage{pdfpages}
\usepackage{supertabular}
\usepackage[table]{xcolor}
\usepackage{soul}
\usepackage{multicol}
\usepackage{paralist}
\usepackage[para]{footmisc}
\usepackage{graphicx}
\usepackage{mdframed}
\usepackage{makecell}

\usepackage{bibunits}

\bibliographyunit[\subsection]
\bibliography*{russo}

\newcommand{\me}[1]{\fcolorbox{black}{white}{\underline{\textbf{#1}}}}
\newcommand{\student}[1]{\fcolorbox{black}{lightgray!40}{#1}}

%% Removes the header "References"
\makeatletter
\renewenvironment{thebibliography}[1]
     {\list{\@biblabel{\@arabic\c@enumiv}}%
           {\settowidth\labelwidth{\@biblabel{#1}}%
            \leftmargin\labelwidth
            \advance\leftmargin\labelsep
            \@openbib@code
            \usecounter{enumiv}%
            \let\p@enumiv\@empty
            \renewcommand\theenumiv{\@arabic\c@enumiv}}%
      \sloppy
      \clubpenalty4000
      \@clubpenalty \clubpenalty
      \widowpenalty4000%
      \sfcode`\.\@m}
     {\def\@noitemerr
       {\@latex@warning{Empty `thebibliography' environment}}%
      \endlist}
\makeatother


%% Compact lists
\newenvironment{CompactItemize}%
  {\begin{list}{$\blacktriangleright$}%
   {\leftmargin=16pt \itemsep=2pt \topsep=2pt
     \parsep=0pt \partopsep=0pt}}%
  {\end{list}}
\renewcommand{\labelitemi}{}


\newif\ifimportantpub
\importantpubfalse

%% To highlight publications
\definecolor{lightblue}{RGB}{196,209,199}
\definecolor{subcolor}{RGB}{255,178,102}
\definecolor{ratecolor}{RGB}{255,102,102}


\setlength{\fboxrule}{0pt} %% No borders

\ifimportantpub
  \newcommand{\important}[1]{\fcolorbox{black}{lightblue}{#1}}
\else
  \newcommand{\important}[1]{{#1}}
\fi

\newcommand{\sub}[1]{}
\newcommand{\rate}[1]{}

\begin{document}
\onecolumn

\kawapplication{Secure AI Agents by Construction: Information Flow Control and Differential Privacy for LLM Tool Use}{A. Russo
  (19780228--7875), \emph{PhD awarded on February 2009}}

\section{1.Bibliometric information}

\begin{CompactItemize}

\item \textbf{Peer-reviewed articles}: 78 
\item \textbf{Total citations} (Google scholar): 3258
\item \textbf{h-index} (Google scholar): 31
\item \textbf{i-index} (Google scholar): 55

My publication record spreads across several venues from different research
communities: programming languages, security, and systems.
%
Taking the conference ranking CORE 2026 \footnote{\url{http://portal.core.edu.au/}},
I have 7 publications in \emph{flagship} A+ conferences: 1 IEEE S\&P, 3 x ACM
SIGSAC CCS, 1 ACM SIGPLAN POPL (distinguished work), and 2 x USENIX OSDI.
%
I also have 15 publications in \emph{top-venues} A conferences: 3 x ACM SIGPLAN
ICFP (one publication got the most influencial paper award in 2022), 1 x ACM
SIGPLAN OOPSLA, 1 x IEEE ICST, 7 x IEEE CSF, 2 x ESORICS, and 1 USENIX HotOS,
as well as prestigious journals like JFP (Cambridge Press), JCS (IOS Press),
and TOPLAS (ACM).

\end{CompactItemize}

\section{2. Selection of research output}
My name is in \me{bold}, while student contributions are in \student{gray}. For
all publications listed below, I significantly contributed to the writing of
the paper. 
%
\begin{enumerate}[1]
  %% DP 
  \item \me{A. Russo}, E. Lobo-Vesga, and M. Gaboardi: \emph{Accuracy for
    Differentially Private Quotients by Fractional Uncertainties} in Proc. of
    ACM SIGSAC Computer and Communication Security (\textbf{CCS}), 2025.
  %
    I conceived the seminal idea and did initial experiments. 
  %
    Quantifying accuracy of DP mechanisms is critical for trustworthy and
    useful AI data releases.

\item E. Lobo-Vesga, \me{A. Russo}, and M. Gaboardi: \emph{Sensitivity by
  Parametricity} in Proc. of ACM SIGPLAN Object-Oriented Programming, Systems,
  Languages \& Applications (\textbf{OOPSLA}), 2024.
  %
  I conceived the seminal idea and initial prototype implementation.
  %
  Automatically deriving sensitivity bounds via type theory helps into 
  generating correct DP queries in AI pipelines. 

\item \student{E. Lobo-Vesga}, \me{A. Russo}, M. Gaboardi: \emph{A Programming
  Language for Data Privacy with Accuracy Estimations} in Journal of ACM
  Transactions on Programming Languages and Systems (\textbf{TOPLAS}), 2021.
  %
  I contributed 
  to the design and implementation of several new features of the IEEE S\&P 2020 paper.  %
  %
  This language-based DP approach underpins the component 
  for safe AI data releases.

  \item \student{E. Lobo-Vesga}, \me{A. Russo}, M. Gaboardi:
    \emph{A Programming Framework for Differential Privacy with Accuracy Concentration Bounds} in Proc. of 41th IEEE Symposium on Security and
    Privacy (\textbf{S\&P}), 2020.
  %
  I had the initial key technical idea which then developed in this work.
  %
    The DPella framework for accuracy-aware DP query
    generation directly informs the DP component of this proposal.

  %% IFC 

\item \student{M. Vassena}, \me{A. Russo}, D. Garg,
  \student{V. Rajani}, and D. Stefan: \emph{From Fine- to Coarse-Grained Dynamic Information Flow Control and Back} in Proc. of ACM SIGPLAN Symposium
  on Principles of Programming Languages (\textbf{POPL})
  (\textbf{distinguished paper}), 2019.
  %
  I came up with the idea and techniques to complete our mechanized proofs.
  %
  This work provides theoretical foundations for dynamic IFC. 



\item \student{T. Schmitz}, \student{M. Algehed}, C. Flanagan,
  and \me{A. Russo}: \emph{Faceted Secure Multi-Execution} in Proc. of ACM
  SIGSAC Computer and Communications Security (\textbf{CCS}), 2018.
  %
  I did the initial implementations as well as a novel proof technique (candidate labels).
  %
  This work could be applied to separate untrusted data read by AI agents
  from trusted descriptions, as suggested by dual-LLMs approach.

%% Old IFC

\item \student{D. Stefan}, \student{E.Z. Yang}, \student{P. Marchenko}, \me{A. Russo}, D. Herman, B. Karp, and D. Mazi\`{e}res: \emph{Protecting Users by Confining JavaScript with COWL} in Proc.
  of 11th USENIX Symposium on Operating Systems Design and Implementation
  (\textbf{OSDI}), 2014.
  %
  I contributed to the main design and use cases. 
  %
  COWL demonstrates practical IFC-based confinement of
  untrusted code. 

\item \student{D. Stefan}, \me{A. Russo}, \student{P. Buiras}, \student{A.
  Levy}, J.C. Mitchell, and D. Mazi\`{e}res: \emph{Addressing Covert
  Termination and Timing Channels in Concurrent Information Flow Systems} in
  Proc. of 17th ACM SIGPLAN International Conference on Functional Programming
  (\textbf{ICFP}), 2012 (\textbf{most influential paper award}).
  %
  I introduced the key ideas. 
  %
  Enforcement of IFC for concurrent systems can be applied to AI pipelines. 

\item \student{D.B. Giffin}, \student{A. Levy}, \student{D. Stefan},
  \student{D. Terei}, D. Mazi\`{e}res, J.C. Mitchell, and
  \me{A. Russo}: \emph{Hails: Protecting Data Privacy in Untrusted Web Applications} in Proc. of 10th USENIX Symposium on Operating Systems Design
  and Implementation (\textbf{OSDI}), 2012.
  %
  I contributed to the design of the system. 
  %
  This work provides a blueprint for confining untrusted MCP tools. 

\item T. Tsai, \me{A. Russo}, and J. Hughes: \emph{A Library for Secure
  Multi-threaded Information Flow in Haskell} in Proc. of 20th IEEE Computer
  Security Foundations Symposium (\textbf{CSF}), 2007.
  %
  I developed some of the IFC tracking  techniques.
  %
  This work directly informs the proposal's use of Arrow-based DSLs to enforce
  control-flow integrity in AI agent planning pipelines.

\end{enumerate}

\section{3. Relevant peer-reviewed research outputs from 2018--2026}

All publications below are peer-reviewed conference papers or journal articles.

\begin{enumerate}[1]

  %% 2025
  \item \me{A. Russo}, E. Lobo-Vesga, and M. Gaboardi: \emph{Accuracy for
    Differentially Private Quotients by Fractional Uncertainties} in Proc. of
    ACM SIGSAC Computer and Communication Security (\textbf{CCS}), 2025.

  %% 2024
\item E. Lobo-Vesga, \me{A. Russo}, and M. Gaboardi: \emph{Sensitivity by
  Parametricity} in Proc. of ACM SIGPLAN Object-Oriented Programming, Systems,
  Languages \& Applications (\textbf{OOPSLA}), 2024.

\item M.E. Cambronero, M.A. Mart\'inez-Pietro, L.F. Llana D\'iaz, and
  \me{A. Russo}: \emph{Towards a GDPR-compliant cloud architecture with data
  privacy controlled through sticky policies} in PeerJ Computer Science,
  Vol.~10, 2024.

  %% 2023
\item \student{A. Sarkar}, \student{R. Krook}, \me{A. Russo}, and K. Claessen:
  \emph{HasTEE: Programming Trusted Execution Environments with Haskell} in
  Proc. of the 16th ACM SIGPLAN International Symposium on Haskell
  (\textbf{Haskell}), 2023.

\item \student{M. Pettersson}, \student{J. Ljung Ekeroth}, and \me{A. Russo}:
  \emph{Calculating Function Sensitivity for Synthetic Data Algorithms} in
  Proc. of the 35th Symposium on Implementation and Application of Functional
  Languages (\textbf{IFL}), 2023.

\item \student{M. Vassena}, \me{A. Russo}, D. Garg, \student{V. Rajani}, and
  D. Stefan: \emph{From Fine- to Coarse-Grained Dynamic Information Flow
  Control and Back} in Foundations and Trends in Programming Languages,
  Vol.~8, No.~1, pp.~1--117, 2023.

  %% 2021
\item \student{E. Lobo-Vesga}, \me{A. Russo}, M. Gaboardi: \emph{A Programming
  Language for Data Privacy with Accuracy Estimations} in Journal of ACM
  Transactions on Programming Languages and Systems (\textbf{TOPLAS}), 2021.

  %% 2020
\item \student{E. Lobo-Vesga}, \me{A. Russo}, M. Gaboardi:
  \emph{A Programming Framework for Differential Privacy with Accuracy
  Concentration Bounds} in Proc. of 41th IEEE Symposium on Security and
  Privacy (\textbf{S\&P}), 2020.

\item \student{C. Tom\'e Corti\~nas}, \student{M. Vassena}, and \me{A. Russo}:
  \emph{Securing Asynchronous Exceptions} in Proc. of the 33rd IEEE Computer
  Security Foundations Symposium (\textbf{CSF}), 2020.

\item \student{N. Valliappan}, \student{R. Krook}, \me{A. Russo}, and K. Claessen:
  \emph{Towards Secure IoT Programming in Haskell} in Proc. of the 13th ACM
  SIGPLAN International Symposium on Haskell (\textbf{Haskell}), 2020.

  %% 2019
\item \student{M. Vassena}, \me{A. Russo}, D. Garg,
  \student{V. Rajani}, and D. Stefan: \emph{From Fine- to Coarse-Grained
  Dynamic Information Flow Control and Back} in Proc. of ACM SIGPLAN Symposium
  on Principles of Programming Languages (\textbf{POPL})
  (\textbf{distinguished paper}), 2019.

\item \student{M. Algehed}, \me{A. Russo}, and C. Flanagan:
  \emph{Optimising Faceted Secure Multi-Execution} in Proc. of the 32nd IEEE
  Computer Security Foundations Symposium (\textbf{CSF}), 2019.

  %% 2018
\item \student{M. Vassena}, \me{A. Russo}, and \student{P. Buiras}: \emph{MAC:
  A verified static information-flow control library} in the Journal of Logical
  and Algebraic Methods in Programming. Vol.~95, 2018.

\item \student{T. Schmitz}, \student{M. Algehed}, C. Flanagan,
  and \me{A. Russo}: \emph{Faceted Secure Multi-Execution} in Proc. of ACM
  SIGSAC Computer and Communications Security (\textbf{CCS}), 2018.

\end{enumerate}

\section{4. Relevant non peer-reviewed research outputs from 2018--2026}

Non-applicable.

\label{lastpage}



\end{document}
